\subsection{GPT-5.3-Codex mit strukturiertem Prompt und Rollendefinition}
Im dritten Experiment (\ref{tab:GPT-5.3-Codex5.3 strukturierter prompt und Rollendefinition}) wurde GPT-5.3-Codex mittels Role Prompting angewiesen, die Rolle eines erfahrenen Softwarearchitekten und Full-Stack-Entwicklers einzunehmen. Es wird damit untersucht, ob die explizite Rollendefinition die Qualität und Vollständigkeit der generierten Lösung weiter verbessert.

Die Evaluation des mit GPT-5.3-Codex  unter Verwendung eines strukturierten Prompts mit Rollendefinition generierten Systems in (\ref{tab:GPT-5.3-Codex5.3 strukturierter prompt und Rollendefinition}) zeigte mehrere Einschränkungen bei der Umsetzung komplexerer Geschäftsprozesse.

\newpage
\begin{table}[H]
\centering
\small
\caption{Ergebnisse von GPT-5.3-Codex mit strukturiertem Prompt und Rollendefinition}
\label{tab:GPT-5.3-Codex5.3 strukturierter prompt und Rollendefinition}
\renewcommand{\arraystretch}{1.15}

\begin{tabularx}{\textwidth}{@{}>{\RaggedRight\arraybackslash}p{5.0cm} >{\centering\arraybackslash}p{1.4cm} >{\RaggedRight\arraybackslash}X@{}}
\toprule
\textbf{Evaluationskriterium} & \textbf{Wert} & \textbf{Beobachtung} \\
\midrule
\multicolumn{3}{@{}l@{}}{\textbf{Funktionale Kriterien}} \\
\addlinespace[0.2em]
Allgemeine Funktionalität & 2 &
Grundlegende Benutzeroberfläche, Formulare und Authentifizierung vorhanden. Erweiterte Funktionen wie Filterung, Pagination sowie vollständige CRUD-Operationen sind vorhanden. \\
Geschäftsobjekte und CRUD-Funktionalität & 2 &
\textit{Customers, Orders, Rooms, Tasks und Invoices} wurden umgesetzt. Mehrere Module besitzen vollständige CRUD-Funktionalität. \\
Erweiterte Domänenfunktionen & 2 &
\textit{\textit{Order Items, Occupancy Items, Services, Kommentare und Room Locations}} wurden vollständig implementiert. \\
Benutzer- und Rechteverwaltung & 2 &
Benutzer- und Rollenmodelle sind nicht vorhanden. \\
\midrule
\multicolumn{3}{@{}l@{}}{\textbf{Technische Kriterien}} \\
\addlinespace[0.2em]
Modell-Validierungen und Geschäftslogik & 1 &
Modell-Validierungen, Geschäftsregeln und technische Qualitätssicherungsmechanismen wurden teilweise umgesetzt. \\
Unit-Tests & 2 &
Es wurden automatisierte Unit-Tests generiert. \\
Codequalität und Wartbarkeit & 2 &
Die technische Struktur enthält nachweisbaren Mechanismen zur Qualitätssicherung oder Testbarkeit. \\
\midrule
\multicolumn{3}{@{}l@{}}{\textbf{Integrationskriterien}} \\
\addlinespace[0.2em]
Frontend-Backend-Integration & 1 &
Grundlegende Interaktionen zwischen Benutzeroberfläche und Backend sind teilweise vorhanden. Komplexe Workflows wurden jedoch nicht vollständig umgesetzt. \\
End-to-End-Workflows & 1 &
Benutzerabläufe oder automatisierten Integrationstests teilweise vorhanden. \\
\midrule
\multicolumn{3}{@{}l@{}}{\textbf{Gesamtergebnis}} \\
\addlinespace[0.2em]
Erreichte Gesamtpunktzahl & 15 / 18 &
Das generierte System erreicht 15 von 18 möglichen Bewertungspunkten und erfüllt damit etwa 83,3\% der definierten Evaluationskriterien. \\
\bottomrule
\end{tabularx}
\end{table}

Die \ref{tab:GPT-5.3-Codex5.3 strukturierter prompt und Rollendefinition} zeigt, dass insbesondere die Funktion zur Erstellung neuer Rechnungen nicht vollständig ausgeführt werden konnte.
Obwohl die entsprechende Benutzeroberfläche zur Anlage einer Rechnung vorhanden war, führte der Speichervorgang nicht zum erfolgreichen Anlegen eines Datensatzes.
Die Funktionalität war somit lediglich teilweise implementiert.
Darüber hinaus traten Probleme beim Anlegen neuer Vorgänge auf.
Der vorgesehene Workflow umfasst mehrere aufeinanderfolgende Schritte, darunter die Auswahl von Räumen, die Buchung zusätzlicher Leistungen, die Zuordnung eines Kunden sowie eine abschließende Zusammenfassung der erfassten Daten.
Während die einzelnen Zwischenschritte grundsätzlich durchlaufen werden konnten, war ein erfolgreicher Abschluss des Prozesses nicht möglich. Insbesondere im letzten Schritt der Zusammenfassung konnte der Vorgang nicht gespeichert werden, sodass keine vollständige Buchung erzeugt wurde.
Die Ergebnisse deuten darauf hin, dass das Modell zwar in der Lage war, die grundlegende Struktur mehrstufiger Geschäftsprozesse zu generieren, jedoch Schwierigkeiten bei der konsistenten Implementierung komplexer Abläufe und deren abschließender Persistierung aufweist. Insbesondere die Verarbeitung mehrstufiger Workflows mit mehreren abhängigen Entitäten konnte nicht vollständig umgesetzt werden.
Die Löschoperation eines Kundendatensatzes führte während der Ausführung zu einem \texttt{ProtectedError}. Die anschließende Analyse ergab jedoch, dass es sich hierbei nicht um einen Implementierungsfehler der Geschäftslogik handelte. Vielmehr war die zugrunde liegende Funktionalität korrekt umgesetzt. Durch die Verwendung der Django-Konfiguration \texttt{on\_delete=models.PROTECT} wurde sichergestellt, dass Kundendatensätze nicht gelöscht werden können, solange sie noch von anderen Systemkomponenten referenziert werden. Im vorliegenden Fall bestand eine Referenz zwischen dem Kunden und einem Datensatz des Modells \texttt{OccupancyLog}. Die Verhinderung der Löschoperation entspricht den Anforderungen zur Wahrung der referenziellen Integrität und dient der Vermeidung inkonsistenter Datenbestände.
Verbesserungspotenzial zeigte sich jedoch in der Behandlung dieser Situation auf Ebene der Benutzerschnittstelle. Anstatt dem Benutzer eine verständliche und fachlich formulierte Meldung anzuzeigen, wurde die technische Ausnahme des Frameworks direkt ausgegeben. Aus Sicht der Benutzerfreundlichkeit wäre eine anwendungsbezogene Fehlermeldung vorzuziehen, die den Grund der fehlgeschlagenen Löschoperation erläutert. Der Vorfall stellt somit keinen funktionalen Fehler der KI-generierten Implementierung dar, sondern eine Schwäche in der Fehlerbehandlung und Benutzerführung.

Bei der Implementierung der generierten Systems wurde ein Fehler im Modul zur Rechnungsverwaltung festgestellt. Nach der Erstellung einer neuen Rechnung und dem Auslösen des Speichervorgangs trat während des Renderns der Rechnungsübersicht ein Laufzeitfehler auf. Django meldete dabei die Fehlermeldung \textit{``Reverse for 'invoice\_update' not found''}. Diese Meldung weist darauf hin, dass im Frontend-Template eine URL mit dem Namen \texttt{invoice\_update} referenziert wurde, für die jedoch keine entsprechende Definition im URL-Routing des Systems vorhanden war.
Die Ursache des Problems lag in einer Inkonsistenz zwischen den verschiedenen Schichten der Anwendung. Während das Template der Rechnungsübersicht davon ausging, dass eine Funktion zur Bearbeitung bestehender Rechnungen existiert, war entweder die zugehörige URL-Konfiguration oder die entsprechende View nicht implementiert. Dadurch konnte Django die angeforderte URL beim Rendern der Seite nicht auflösen und brach die Verarbeitung mit einer Ausnahme ab.
Aus Sicht der Softwarearchitektur handelt es sich um einen Integrationsfehler zwischen Präsentationsschicht, Routing-Konfiguration und Geschäftslogik. Obwohl die einzelnen Artefakte des Systems syntaktisch korrekt generiert wurden, wurde ihre gegenseitige Konsistenz nicht vollständig sichergestellt. Der Fehler verdeutlicht eine typische Schwäche generativer KI-Modelle bei der Entwicklung mehrschichtiger Anwendungen: Das Modell erzeugt einzelne Komponenten, ohne sämtliche Abhängigkeiten zwischen Templates, Views und URL-Konfigurationen zuverlässig zu berücksichtigen.
Die Auswirkungen dieses Fehlers gingen über die reine Anzeige einer Fehlermeldung hinaus. Da die Rechnungsübersicht nicht mehr korrekt gerendert werden konnte, wurde der Arbeitsablauf im Rechnungsmodul unterbrochen. Benutzer konnten die erzeugten Rechnungen nicht wie vorgesehen verwalten oder bearbeiten.

Dies reduzierte die Nutzbarkeit des Moduls erheblich und machte eine manuelle Fehleranalyse sowie die nachträgliche Ergänzung der fehlenden Routing- oder View-Komponenten erforderlich.
Der Fall zeigt, dass die Qualität des generierten Software nicht ausschließlich anhand der syntaktischen Korrektheit des erzeugten Codes bewertet werden kann. Vielmehr ist eine systematische Überprüfung der Konsistenz zwischen den verschiedenen Anwendungsschichten erforderlich. Insbesondere bei Webanwendungen müssen Templates, URL-Konfigurationen und Views gemeinsam validiert werden, um Integrationsfehler dieser Art frühzeitig zu erkennen und zu beheben.
